
State-of-the-art language models that achieve high scores on one trustworthiness benchmark systematically excel on others, including benchmarks measuring seemingly unrelated capabilities such as mathematical reasoning and instruction following. Across the Open LLM Leaderboard's 4,561 models, benchmark scores correlate at $\rho$ = 0.80--0.87. This observation raises questions about a foundational assumption in LLM evaluation: that trustworthiness dimensions are independent constructs requiring separate measurement.

The standard approach treats each dimension as a distinct target. Specialized benchmarks---TruthfulQA for factual accuracy, MMLU for knowledge, AdvGLUE for robustness---are developed and deployed independently. Practitioners evaluate models dimension-by-dimension, constructing scorecards that imply orthogonal capabilities. Yet the observed cross-correlation suggests this framing may not reflect how trustworthiness manifests in language models.

Prior work has typically attributed shared variance to model scale and dismissed residual correlation as noise. This study instead investigates whether the residual structure contains meaningful signal. Drawing on psychometric methodology, where similar positive manifold patterns led to factor-analytic discoveries, we examine whether an analogous latent construct underlies trustworthy LLM behavior. We term this hypothesized factor Generalized Representational Coherence (GRC), positing that it may arise from representation stability across semantically equivalent inputs.

This work makes the following contributions:

\begin{enumerate}
\item \textbf{Existence of residual factor:} Through meta-analysis (N = 4,561 models), a dominant principal component (PC1, $\lambda _1$ = 2.277, explaining 60\% of residual variance) persists after controlling for log(parameters) and release date, exceeding permutation null thresholds (p = 0.001).
\end{enumerate}

\begin{enumerate}
\item \textbf{Mechanistic evidence (proof-of-concept):} PC1 correlates positively with a synthetic Behavioral Stability Index ($\rho$ = 0.405, p < $10^{-179})$. This correlation is suggestive but requires validation with real behavioral measurements.
\end{enumerate}

\begin{enumerate}
\item \textbf{Quasi-intervention evidence:} Instruction-tuning increases both BSI (Cohen's d = 1.87) and PC1 scores (d = 1.99) within 16 matched base/instruct model pairs.
\end{enumerate}

\begin{enumerate}
\item \textbf{Prospective validity (simulated):} Frozen PC1 weights generalize to 5 of 6 simulated holdout benchmarks (loadings $\geq$ 0.3). True prospective validation awaits independent benchmark data.
\end{enumerate}
